\section{Model}
\label{sec:model}

The model is a standard real business-cycle economy with two ways of producing: a formal sector that uses capital and registered labor, and an informal sector made of own-account work and small units that uses labor only. It is kept as close as possible to the textbook structure so that the identification results are not driven by ad hoc features. This section derives every equilibrium condition from the optimization problems of households and firms, explains what each condition means, and shows how the informal share of output follows from them. The stationary system that is solved numerically, the steady-state algorithm and the linear solution are collected in Appendix~\ref{app:model}.

\subsection{Overview and notation}

Time is quarterly. The economy has a representative household, competitive formal firms, informal production units owned and run by the household, and a government. Table~\ref{tab:notation} lists the main symbols. Capital letters denote levels that grow with the economy; the corresponding lowercase letters, used from Section~\ref{sec:steady} on, denote the same variables divided by the trend.

\begin{table}[H]
\centering
\caption{Main notation.}
\label{tab:notation}
\footnotesize
\begin{tabularx}{\textwidth}{l X}
\toprule
\textbf{Symbol} & \textbf{Meaning} \\
\midrule
$C_t, I_t, G_t$ & Private consumption, investment, government consumption \\
$K_t$ & Capital stock available in quarter $t$ (all capital is used by formal firms) \\
$N_{f,t}, N_{i,t}$ & Formal and informal work \\
$s_t = N_{i,t}/(N_{f,t}+N_{i,t})$ & Labor-informality rate \\
$Y_{f,t}, Y_{i,t}, Y_t$ & Formal, informal and total output, $Y_t = Y_{f,t}+Y_{i,t}$ \\
$\theta_t = Y_{i,t}/Y_t$ & Informal share of output; $\theta$ without subscript is its steady state \\
$w_{f,t}, r_t$ & Formal wage and rental rate of capital \\
$\tau_t$ & Formality wedge on formal labor income (taxes, contributions, compliance costs) \\
$X_t$ & Labor-augmenting technology, common to both sectors, with a unit root \\
$R_t$ & Informal income per worker relative to the gross formal wage \\
$\alpha$ & Capital share of formal output \\
$\eta$ & Elasticity of informal output to informal work ($0<\eta<1$) \\
$\varphi$ & Inverse Frisch elasticity of labor supply \\
$\mu_t, z^g_t, \nu_t, \psi_t, \omega_t, x_t, z^r_t$ & Shocks: trend growth, government, investment efficiency, labor supply, wedge, absorption, informal productivity \\
\bottomrule
\end{tabularx}
\end{table}

\subsection{Household}
\label{sec:household}

\paragraph{Preferences.} A representative household chooses consumption, investment, next quarter's capital and the hours supplied to each sector to maximize expected lifetime utility
\begin{equation}
E_0 \sum_{t=0}^{\infty}\beta^t\Big[\ln C_t - \psi_t\,\frac{\chi_f N_{f,t}^{1+\varphi} + \chi_i N_{i,t}^{1+\varphi}}{1+\varphi}\Big].
\label{eq:utility}
\end{equation}
The first term is the satisfaction from consumption; the logarithm makes the marginal value of an extra peso fall as consumption rises. The second term is the disutility of work. Each type of work has its own convex cost: the hundredth hour in the informal sector is more costly than the first, and the same holds for formal work. This is the simplest specification in which the household supplies both types of work at the same time, as Colombian households do, instead of moving everybody to whichever sector pays slightly more. The weights $\chi_f$ and $\chi_i$ capture fixed differences in the attractiveness of the two sectors (access, flexibility, preferences) and are calibrated so that the model reproduces the mean informality rate. The shock $\psi_t$ shifts the disutility of all work (labor supply), $\beta$ is the discount factor and $1/\varphi$ is the Frisch elasticity of labor supply.

\paragraph{Budget and capital accumulation.} Each quarter the household receives formal wages net of the wedge, the whole output of the informal units it runs, the rental income of its capital, lump-sum transfers $T_t$ and an aggregate net-external-absorption transfer $B_t$, which is taken as given in the household optimization. It spends on consumption and investment:
\begin{equation}
C_t + I_t = (1-\tau_t)\,w_{f,t}N_{f,t} + Y_{i,t} + r_t K_t + T_t + B_t .
\label{eq:budget}
\end{equation}
Investment becomes capital with an efficiency that varies over time,
\begin{equation}
K_{t+1} = (1-\delta)K_t + \nu_t I_t ,
\label{eq:capital}
\end{equation}
where $\delta$ is depreciation and $\nu_t$ is a shock to the marginal efficiency of investment: a peso invested produces more or less new capital depending on financial conditions or the relative price of capital goods. The formality wedge is
\begin{equation}
1-\tau_t = (1-\bar\tau)\,e^{-\omega_t},
\label{eq:wedge}
\end{equation}
so that $\bar\tau$ is its steady-state level and an increase in $\omega_t$ makes formal work less attractive. The wedge summarizes everything that separates what a formal job costs from what the worker keeps: social-security contributions, payroll taxes and the costs of registration and compliance.

\paragraph{First-order conditions.} Let $\lambda_t$ be the Lagrange multiplier of the budget constraint (the marginal utility of a peso) and $\lambda_tq_t$ that of the capital accumulation equation (the value of a unit of installed capital). The Lagrangian is
\begin{multline}
\mathcal{L} = E_0\sum_{t=0}^{\infty}\beta^t\Big\{\ln C_t - \psi_t\frac{\chi_fN_{f,t}^{1+\varphi}+\chi_iN_{i,t}^{1+\varphi}}{1+\varphi}
+ \lambda_t\big[(1-\tau_t)w_{f,t}N_{f,t} + Y_{i,t} + r_tK_t + T_t + B_t - C_t - I_t\big] \\
+ \lambda_tq_t\big[(1-\delta)K_t + \nu_tI_t - K_{t+1}\big]\Big\}.
\end{multline}
Differentiating with respect to each choice and setting the derivative to zero gives five conditions.
\begin{align}
C_t:&\quad \lambda_t = \frac{1}{C_t}, \label{eq:foc-c}\\
I_t:&\quad q_t = \frac{1}{\nu_t}, \label{eq:foc-i}\\
K_{t+1}:&\quad \lambda_tq_t = \beta E_t\big[\lambda_{t+1}\big(r_{t+1} + q_{t+1}(1-\delta)\big)\big], \label{eq:foc-k}\\
N_{f,t}:&\quad \psi_t\chi_fN_{f,t}^{\varphi} = \lambda_t(1-\tau_t)w_{f,t}, \label{eq:foc-nf}\\
N_{i,t}:&\quad \psi_t\chi_iN_{i,t}^{\varphi} = \lambda_t\,\frac{\partial Y_{i,t}}{\partial N_{i,t}} . \label{eq:foc-ni}
\end{align}
Each has a direct reading. \eqref{eq:foc-c}: the value of a peso is the extra utility of consuming it. \eqref{eq:foc-i}: a unit of capital costs $1/\nu_t$ pesos of investment, so that is its price. \eqref{eq:foc-k}: the household gives up a unit of capital today, worth $\lambda_tq_t$, only if it expects to recover the same value tomorrow from the rent $r_{t+1}$ plus the undepreciated capital, discounted by $\beta$. Substituting \eqref{eq:foc-c} and \eqref{eq:foc-i} gives the Euler equation
\begin{equation}
\frac{1}{\nu_t C_t} = \beta E_t\Big[\frac{1}{C_{t+1}}\Big(r_{t+1} + \frac{1-\delta}{\nu_{t+1}}\Big)\Big].
\label{eq:euler}
\end{equation}
The two labor conditions say that the household works in each sector until the utility cost of the last hour equals the utility value of what that hour produces. Using \eqref{eq:foc-c} to replace $\lambda_t$,
\begin{equation}
\psi_t\chi_f N_{f,t}^{\varphi} C_t = (1-\tau_t)\,w_{f,t}, \qquad \psi_t\chi_i N_{i,t}^{\varphi} C_t = \frac{\partial Y_{i,t}}{\partial N_{i,t}} .
\label{eq:foc-labor}
\end{equation}
The left-hand side is the marginal rate of substitution between leisure and consumption: how many pesos of consumption the household asks for to work one more hour. The right-hand side is what that hour pays: the formal wage after the wedge in one case, the marginal product of informal work in the other. A higher wedge lowers the right-hand side of the formal condition, so the household shifts hours towards informal work until both conditions hold again.

\subsection{Formal firms}

Formal firms are competitive, rent capital from the household and hire formal work to produce with a Cobb--Douglas technology,
\begin{equation}
Y_{f,t} = K_t^{\alpha}\big(X_tN_{f,t}\big)^{1-\alpha}.
\label{eq:technology}
\end{equation}
A firm maximizes profits $Y_{f,t} - w_{f,t}N_{f,t} - r_tK_t$, which gives
\begin{equation}
r_t = \alpha\,\frac{Y_{f,t}}{K_t}, \qquad w_{f,t} = (1-\alpha)\,\frac{Y_{f,t}}{N_{f,t}} .
\label{eq:foc-firm}
\end{equation}
Each factor is paid its marginal product. Because the technology has constant returns, the payments exhaust output: a share $\alpha$ goes to capital and a share $1-\alpha$ to formal labor. This is the property used to estimate $\alpha$ (Section~\ref{sec:alpha}), and it is what makes the formal wage measured in the survey a fraction $1-\alpha$ of formal value added per worker.

Labor-augmenting technology has a unit root with stochastic growth \citep{AguiarGopinath2007},
\begin{equation}
\ln X_t - \ln X_{t-1} = \gamma + \mu_t ,
\label{eq:trend}
\end{equation}
so a shock $\mu_t$ shifts the level of productivity permanently. This is the standard way to represent the persistent swings in growth of emerging economies, and it means that output, consumption, investment and capital share a common stochastic trend.

\subsection{Informal production}

Following the shadow-economy literature \citep{IhrigMoe2004, ElginOztunali2012}, informal production uses labor only, with decreasing returns:
\begin{equation}
Y_{i,t} = A_i\,e^{z^r_t}\,X_t\,N_{i,t}^{\eta}, \qquad 0<\eta<1 .
\label{eq:informal-technology}
\end{equation}
$A_i$ is the level of informal productivity, $z^r_t$ a shock to it, and $X_t$ enters so that informal output grows with the economy. Decreasing returns ($\eta<1$) mean that each additional informal worker adds less than the previous one: informal activities compete for the same customers and spaces, and they cannot be scaled up without capital.

\paragraph{Average and marginal product.} The informal unit is run by the household, so there is no separate informal wage. Two quantities matter and must be kept apart:
\begin{equation}
\underbrace{\frac{Y_{i,t}}{N_{i,t}}}_{\text{average product: income per informal worker}}
\qquad\text{and}\qquad
\underbrace{\frac{\partial Y_{i,t}}{\partial N_{i,t}} = \eta\,\frac{Y_{i,t}}{N_{i,t}}}_{\text{marginal product: return of one more hour}} .
\label{eq:avg-marginal}
\end{equation}
The marginal product enters the household's choice of hours in \eqref{eq:foc-labor}. The average product is the latent value added per worker that the measurement exercise seeks to approximate with survey earnings. It is not automatically equal to every respondent's reported wage or profit.

\paragraph{Why labor only.} The justification is not that informal units use no capital, but what the household survey measures. The income reported by own-account workers and informal employers is mixed income or profit: it already contains the return to the vehicles, tools and premises of their own unit, and the wages they pay are counted in the income of their employees. Informal survey income is therefore a measure of informal value added per worker. The income of formal wage employees, by \eqref{eq:foc-firm}, represents only the labor share of formal value added. The labor-only technology implements these measurement conventions \emph{inside the model}: informal income per worker is defined as informal output per worker, whereas the formal wage is a fraction $1-\alpha$ of formal output per worker. Employees of informal units are the one group whose reported wage does not contain the unit's surplus, and the earnings series is adjusted for them by group (Section~\ref{sec:data}); alternative treatments are reported in the sensitivity analysis.

\subsection{Government, resources and shocks}

Government consumption is exogenous and grows with the trend, $G_t = \bar g\,X_t\,e^{z^g_t}$, and the government budget is balanced with lump-sum transfers, $T_t = \tau_tw_{f,t}N_{f,t} - G_t$. The wedge is therefore modeled as a transfer between agents, not a direct loss of resources; it changes the allocation of work. To account for national-account components outside the domestic three-good model, let $B_t=(\bar x+x_t)Y_t$ denote a reduced-form net transfer financing excess domestic absorption (or, when negative, net resources supplied to the rest of the world). It is taken as given by individual households. This is an accounting closure rather than a structural model of an external sector. Aggregating the household and government budgets and using \eqref{eq:foc-firm} gives
\begin{equation}
Y_t\,\big(1 + \bar x + x_t\big) = C_t + I_t + G_t .
\label{eq:resources}
\end{equation}
The term $\bar x+x_t$ represents omitted absorption and the statistical discrepancy arising from chained volumes. In particular, $Y=C+I+G$ is not a valid identity for Colombia's open economy or for non-additive chained-volume components. The steady-state term $\bar x$ reproduces the sample expenditure ratios; $x_t$ absorbs fluctuations in the omitted components. Because $B_t$ has no intertemporal law of repayment, this closure must not be interpreted as a structural theory of trade, external debt or fiscal financing.

The seven exogenous processes are independent first-order autoregressions,
\begin{equation}
\zeta_t = \rho_\zeta\,\zeta_{t-1} + \sigma_\zeta\,\varepsilon_{\zeta,t},\qquad \varepsilon_{\zeta,t}\sim N(0,1),\qquad \zeta\in\{\mu,\, z^g,\, \ln\nu,\, \ln\psi,\, \omega,\, x,\, z^r\},
\label{eq:shocks}
\end{equation}
with persistence $\rho_\zeta$ and volatility $\sigma_\zeta$. Their roles are distinct: $\mu$ moves the growth of productivity, $z^g$ government spending, $\nu$ the efficiency of investment, $\psi$ the willingness to work, $\omega$ the cost of formality, $x$ absorption and $z^r$ informal productivity. The informal-productivity shock is active only when the earnings observable of Section~\ref{sec:identification} is used; without it nothing in the data separates it from the wedge, and it is switched off in the estimates based on aggregates only.

\paragraph{Equilibrium.} An equilibrium is a set of quantities and prices that satisfy the household conditions \eqref{eq:capital}, \eqref{eq:euler} and \eqref{eq:foc-labor}, the firm conditions \eqref{eq:foc-firm}, the technologies \eqref{eq:technology} and \eqref{eq:informal-technology}, the resource constraint \eqref{eq:resources} and the laws of motion of the shocks. After substituting prices, the model has six endogenous variables ($C$, $Y$, $I$, $N_f$, $N_i$, $K$) and seven shocks, thirteen equations in total (Appendix~\ref{app:model}).

\subsection{The informal share of output}

\begin{proposition}[Informal share of output]
\label{prop:share}
Let $R_t$ be informal income per worker relative to the gross formal wage, $R_t \equiv (Y_{i,t}/N_{i,t})/w_{f,t}$, and $s_t \equiv N_{i,t}/(N_{f,t}+N_{i,t})$ the informality rate. Then
\begin{equation}
\theta_t \equiv \frac{Y_{i,t}}{Y_t} = \frac{(1-\alpha)\,R_t\,o_t}{1+(1-\alpha)\,R_t\,o_t},\qquad o_t\equiv\frac{s_t}{1-s_t}.
\label{eq:share}
\end{equation}
If instead both sectors shared the Cobb--Douglas technology with mobile capital, the factor $1-\alpha$ would disappear and $\theta_t$ would equal the informal share of labor income, $R_to_t/(1+R_to_t)$.
\end{proposition}

\emph{Proof.} Write the ratio of informal to formal output as a product of output per worker and workers:
\[
\frac{Y_{i,t}}{Y_{f,t}} = \frac{Y_{i,t}/N_{i,t}}{Y_{f,t}/N_{f,t}}\cdot\frac{N_{i,t}}{N_{f,t}}
= \frac{R_t\,w_{f,t}}{w_{f,t}/(1-\alpha)}\cdot o_t = (1-\alpha)\,R_t\,o_t ,
\]
where $o_t$ denotes the odds of informal employment and the second step uses the definition of $R_t$ and $w_{f,t}=(1-\alpha)Y_{f,t}/N_{f,t}$ from \eqref{eq:foc-firm}, and $N_{i,t}/N_{f,t}=s_t/(1-s_t)$. Then $\theta_t = (Y_{i,t}/Y_{f,t})/(1+Y_{i,t}/Y_{f,t})$. With a common technology both sectors have the same capital per efficiency unit of labor, output is proportional to labor income in each sector, and the factor $1-\alpha$ cancels. $\square$

Proposition~\ref{prop:share} is an accounting identity of the model, and it separates the measurement problem in three parts.
\begin{itemize}
  \item The informality rate $s_t$ is observed directly in the household survey.
  \item The earnings ratio $R_t$ is the object to identify. It is also measured in the survey, as the ratio of mean informal income to the mean wage of formal employees, after removing the wedge.
  \item The capital share $\alpha$ converts the formal wage into formal output per worker. It is estimated from the national accounts.
\end{itemize}
A numerical example shows how much each part matters. At an informality rate of 55\% ($o=1.22$) and $\alpha=0.35$, an earnings ratio of $R=0.2$ gives $\theta=14\%$ and $R=0.6$ gives $\theta=32\%$. With the same $R=0.35$, raising $\alpha$ from 0.35 to 0.50 lowers $\theta$ from 22\% to 18\%, and assuming equal technologies raises it to 30\%. For a fixed $R$ and $s$, the equal-technology expression is larger because it omits the factor $1-\alpha$. For given data it is therefore an upper bound.

\subsection{What moves the informality rate}

Dividing the two labor conditions in \eqref{eq:foc-labor} eliminates consumption and the labor-supply shock, which affect both sectors in the same way, and leaves the relative return of the two types of work:
\begin{equation}
\frac{\chi_i}{\chi_f}\Big(\frac{N_{i,t}}{N_{f,t}}\Big)^{\varphi} = \frac{\eta\,Y_{i,t}/N_{i,t}}{(1-\tau_t)(1-\alpha)\,Y_{f,t}/N_{f,t}} .
\label{eq:relative-return}
\end{equation}
The household moves work towards the informal sector when its marginal return rises relative to the after-wedge formal wage, and the convex costs ($\varphi>0$) limit how far it goes.

\begin{proposition}[Determinants of the informality rate]
\label{prop:rate}
Taking logarithms of \eqref{eq:relative-return}, with $Y_{i,t}/N_{i,t} = A_ie^{z^r_t}X_tN_{i,t}^{\eta-1}$,
\begin{equation}
\varphi\logit s_t + (1-\eta)\ln N_{i,t} + \ln\frac{Y_{f,t}}{X_tN_{f,t}}
= \ln\frac{\chi_f\,\eta A_i}{\chi_i\,(1-\alpha)} - \ln(1-\bar\tau) + z^r_t + \omega_t .
\label{eq:rate-labor-only}
\end{equation}
The right-hand side contains only the wedge and informal productivity. The left-hand side contains, besides the informality rate, the scale of informal work and formal output per efficiency unit of labor, which depends on capital per formal worker. Aggregate shocks to productivity, investment efficiency, labor supply, government spending or absorption therefore move the informality rate only through capital per formal worker and through the scale of informal employment. With equal technologies those two terms disappear and
\begin{equation}
\varphi\,\logit s_t = \ln\frac{\chi_f}{\chi_i} - \ln(1-\bar\tau) + \ln\tilde\rho + z^r_t + \omega_t ,
\label{eq:rate}
\end{equation}
so aggregate shocks do not move the informality rate at all.
\end{proposition}

\emph{Proof.} Take logarithms of \eqref{eq:relative-return}, use $N_{i,t}/N_{f,t} = s_t/(1-s_t)$, $\ln[s_t/(1-s_t)]=\logit s_t$ and \eqref{eq:wedge}, and collect the endogenous terms on the left. With equal technologies $Y_t=K_t^{\alpha}(X_tL_t)^{1-\alpha}$, $L_t = N_{f,t}+\tilde\rho e^{z^r_t}N_{i,t}$, the informal wage is $\tilde\rho e^{z^r_t}w_{f,t}$ and the ratio of wages does not depend on capital or on the scale of either sector. $\square$

Proposition~\ref{prop:rate} has a direct empirical consequence: the variance of the informality rate is attributed mostly to the formality wedge and to informal productivity (Section~\ref{sec:results}). It is also the reason why the informality rate carries little information about $R_t$ beyond its own level.

\paragraph{Equal technologies.} For comparison and as an upper bound, the model is also estimated with a common Cobb--Douglas technology, $Y_t=K_t^{\alpha}(X_tL_t)^{1-\alpha}$ with $L_t=N_{f,t}+\tilde\rho e^{z^r_t}N_{i,t}$, mobile capital and an informal wage $\tilde\rho e^{z^r_t}w_{f,t}$, where $\tilde\rho$ is the relative productivity of informal work. In that case $R_t = \tilde\rho e^{z^r_t}$ and the share of output equals the share of labor income.

\subsection{Trend, steady state and solution}
\label{sec:steady}

\paragraph{Removing the trend.} Output, consumption, investment, government consumption and capital grow with $X_t$, so they have no constant steady state. The model is made stationary by dividing them by the trend: $y_t=Y_t/X_t$, $c_t=C_t/X_t$, and so on, and $k_t=K_t/X_{t-1}$ for capital, which is chosen one quarter in advance. Work and the informality rate are already stationary. The stationary equations are listed in Appendix~\ref{app:model}.

\paragraph{Steady state.} With all shocks at zero the stationary model has a steady state that can be computed in closed form, step by step, for given $\theta$:
\begin{enumerate}[leftmargin=*]
  \item The Euler equation \eqref{eq:euler} with constant consumption growth $e^{\gamma}$ gives the rental rate, $r = e^{\gamma}/\beta - 1 + \delta$. With the calibration of Section~\ref{sec:estimation} it is 3.5\% per quarter.
  \item The firm condition \eqref{eq:foc-firm} gives the capital--output ratio of the formal sector, $K/Y_f = \alpha/r$: 10.0 quarters, or 2.5 years, close to the 2.4--2.8 measured by DANE for the whole economy.
  \item Hours are split according to the mean informality rate: $N_f = \bar N(1-\bar s)$ and $N_i=\bar N\bar s$, with $\bar N$ a normalization.
  \item Formal output follows from the technology, $Y_f = (K/Y_f)^{\alpha/(1-\alpha)}N_f$ in detrended units, and total output from the informal share, $Y = Y_f/(1-\theta)$, so $Y_i=\theta Y$ and $A_i = Y_i/N_i^{\eta}$.
  \item Investment replaces depreciation and keeps capital growing with the trend, $I = (e^{\gamma}-1+\delta)\,\tilde K$ with $\tilde K = K_t/X_t$. The implied ratio $I/Y$ is 14\%, below the 19\% of the data, because the model has no residential or public investment of its own; consumption and government consumption are set at their sample ratios to GDP, and $\bar x$ closes the resource constraint.
  \item Finally, the weights $\chi_f$ and $\chi_i$ are those that make the labor conditions \eqref{eq:foc-labor} hold at the chosen $N_f$ and $N_i$.
\end{enumerate}
The last step is important for identification. Because $A_i$, $\chi_f$ and $\chi_i$ adjust, \emph{any} value of $\theta$ is consistent with the observed informality rate and with all steady-state ratios: the level of $\theta$ is not pinned down by the first moments of the aggregates. The same holds with equal technologies, where $\tilde\rho=\theta(1-\bar s)/[(1-\theta)\bar s]$. Whatever information there is about $\theta$ must come from dynamics or from an additional observable (Section~\ref{sec:identification}).

\paragraph{Solution.} The thirteen stationary equations are linearized around the steady state: each variable is replaced by its log deviation from the steady state, and each equation by its first-order Taylor approximation. The result is a linear system $\bm{A}\,E_t\bm{z}_{t+1} = \bm{B}\,\bm{z}_t$, in which eight variables are predetermined (capital and the seven exogenous processes) and the other five (consumption, output, investment and the two types of work) are choices that can jump when news arrives. The system is solved by the generalized Schur decomposition \citep{Klein2000}, which returns the unique stable solution
\begin{equation}
\bm{s}_{t+1} = \bm{T}\bm{s}_t + \bm{R}\bm{\varepsilon}_{t+1},\qquad \text{jumps}_t = \bm{P}\bm{s}_t ,
\label{eq:solution}
\end{equation}
where $\bm{s}_t$ collects the eight state variables (capital and the seven shocks). A unique stable solution exists when the number of stable eigenvalues equals the number of states (the Blanchard--Kahn condition); this is checked at every parameter draw, and draws that violate it get zero posterior density.
